\begin{abstract}
Retrieval-augmented generation (RAG) treats semantically similar documents as interchangeable evidence, but in law a lexically close precedent may already have been narrowed, superseded, or overruled at the date that matters. We study this failure in a product-facing setting and introduce \emph{Product--Litigation Risk Entailment} (PLRE), a directional task that asks whether an AI product's attributes entail exposure to a legal-risk theory under precedent that is valid at a given query date. Our system, ALERT (Agentic Legal Exposure and Risk Tracking), infers \emph{when} each authority stops being operative: it extracts treatment edges (overrules, supersedes, narrows) from later opinions, turns them into validity intervals, and validates them against gold annotations and a commercial citator (negative-treatment micro-F1 0.81). A deliberately simple retrieval-time gate applies these intervals before generation, and a goal-conditioned agentic loop stops when evidence coverage, not legal correctness, is satisfied. On ALERT-Dataset (1{,}247 U.S.\ AI-litigation cases, 312 product-mapping cases), validity gating raises PLRE F1 from 0.72 to 0.78 and cuts stale citations from 12.7\% to 4.1\% relative to flat similarity, and ALERT improves over rule-based, single-pass LLM, RAG-only, single-agent, and iteration-matched self-refinement baselines on a human-adjudicated subset (macro F1 0.79 vs.\ 0.75 for the strongest baseline). We release the dataset with per-record label provenance, annotation guidelines, and evaluation scripts. ALERT is decision support for qualified reviewers, not legal advice.
\end{abstract}
